\documentclass[11pt,a4paper]{article}


%\setcounter{page}{2}


\begin{document}

%\input{Sections/top}

\section{Course Outline}
The course is on advanced methods in control, especially relevant to Robotics, Industrial, Transport and AI applications. We cover content from Stability certificates to SOS to MPPI to autonomous robot control.

\section{Grading Policy}
Grading criteria:
\begin{itemize}
	\item A: 85-100 points
	\item B: 70-84 points
	\item C: 50-69 points
	\item D: less than 50 points
\end{itemize}

Credit distribution:
\begin{itemize}
	\item 20 points can be awarded for the Assignment 1
	\item 20 points can be awarded for the Assignment 1
	\item 10 points can be awarded for in-class activity
	\item 50 points are awarded for the exam
\end{itemize}

\subsection{Late Submission Policy}
Points might be deducted if the work is submitted after deadline


\section{Course Structure}

We cover the topics:

\begin{enumerate}
	\item Introduction, control theory recap
	\item Nonlinear stability certificates, region of attraction
	\item Koopman, SOS (for stability certificates)
	\item Basic nonlinear control (gravity compensation, CTC, regressor-form adaptive control)
	\item Nonlinear MPC
	\item MPPI
	\item Advanced linear control: time delay, disturbance observer
	\item Advanced linear control: robust MPC
	\item Extended Kalman filter
	\item Trajectory planning: collocation, shooting
	\item RL-based control
	\item Walking robot control: Raibert’s model, Kim’s MPC, WBIC
	\item Walking robot control: Lie group-based observer
	\item UAV control: differential flatness
\end{enumerate}


\section{Course Outline}


\subsection{Knowledge Areas (in terms of application)}

\begin{itemize}
	\item Robotics
	\item Automation
	\item Design (mechanical, electrical)
	\item Control
	\item Data analysis
	\item Computational geometry
\end{itemize}

\subsection{Course Delivery}

The lectures are given every week, followed by practical sessions (labs). During the sessions, students are required to develop their own code based on the knowledge acquired during lectures and the self-study.

\subsection{Prerequisite courses}

\begin{itemize}
	\item Strong prerequisites: Linear Algebra.
	\item Weak prerequisites: Calculus, Control Theory.
	\item Required background knowledge: Python (alternatively Matlab or any other language suitable to work with linear algebra-heavy problems) 
\end{itemize}

\subsection{Expected Learning outcomes}

The course will provide an opportunity for participants to:
\begin{itemize}
	\item Understand and learn to use modern control methods: direct Lyapunov method, sum of squares stability certificates, Koopman representation, etc.
	\item Learn how to use modern control methods for application in Robotics.
	\item Being able to implement a viraety of algorithms covered in te course (see topics).
\end{itemize}


\subsection{Reference Materials}
\begin{itemize}
	\item Annotated slides
	\item Online materials
	\item Educational videos
\end{itemize}

\subsection{Computer Resources}
Students will need to run computer experiments on a laptop and/or on lab computers. 

\subsection{Laboratory Exercises} 
There are a series of labs and electronic handouts prepared for the course.

\subsection{Laboratory Resources}
Students will be required to use and modify a software tool written in Python which run on multiple platforms (Linux, Microsoft Windows, and Mac OS). The tool requires freely available software libraries.

\subsection{Cooperation Policy and Quotations}
We encourage intensive discussion and collaboration in this class. You should feel free to discuss all aspects of the class with classmates and work with them to complete your assignments and project report. However, if you are working together, you must provide details of your contribution and that of others.


\end{document}












